\section{Methodology}

The unit of audit is an exact mathematical assertion, not a paragraph's
plausibility.  Each assertion is recorded with its domain, codomain,
coefficient ring, variance convention, and dependencies.  The four outcomes
used in this paper are literal:

\begin{description}[leftmargin=7.4em,style=nextline]
\item[Validated] A complete proof from the stated hypotheses is given or a
precisely cited theorem is applied after every hypothesis is checked.
\item[Repaired] The printed assertion fails as stated, but a corrected
statement is proved and every downstream use is recomputed.
\item[Disproved] A complete contradiction or counterexample is given to the
exact assertion.  A defect in a proof is not, by itself, a disproof of its
theorem.
\item[Open] No complete proof or counterexample is presently recorded.  This
is not a euphemism for false.
\end{description}

Computations are converted into checkable certificates: displayed matrices,
integer row or column relations, determinant ideals, primitive kernels, or
explicit substitutions.  A software pass is corroboration, not a proof.
Likewise, an algebraically correct matrix calculation does not identify that
matrix with a geometric specialization, attaching, or boundary map; that
identification receives its own statement and proof obligation.

Relevant human literature is searched proactively, not only when a remembered
fact seems doubtful.  Results are cited at the point of use with theorem,
section, or page locators and with their hypotheses intact.  Public working
\LaTeX{} editions such as the maintained SGA English source
\cite{SGAEnglish2026} may be used as searchable reading layers, but the
published SGA volumes \cite{SGA7I1972,SGA7II1973} remain canonical.  Every
reference in the compiled paper is resolvable without access to the editor's
computer.  The citation and tool-disclosure policy follows the Leiden
Declaration \cite{Leiden2026}; neither a model-generated review nor an
internal audit note substitutes for a human mathematical source.

